\section*{Chapter \ref{ch:m2:european-and-japanese-government-bonds} --- European and Japanese Government Bonds}

\begin{solution}{exo:m2:european-and-japanese-government-bonds:1}
Spain $3.633 - 3.18 = 45$, Italy $3.986 - 3.18 = 81$, France $4.00 - 3.18 = 82$
basis points.
\end{solution}

\begin{solution}{exo:m2:european-and-japanese-government-bonds:2}
The spread widens by $30 - (-10) = 40$ basis points: 30 from Italy, 10 from
Germany's fall.
\end{solution}

\begin{solution}{exo:m2:european-and-japanese-government-bonds:3}
\omterm{def:m2:european-and-japanese-government-bonds:ldi}{Gilts} 1\,000, repo 500, cushion 500: leverage 2. After a 20\% fall: \omterm{def:m2:european-and-japanese-government-bonds:ldi}{gilts} 800,
repo 500, cushion 300 (60\% of the original), leverage $800/300 = 2.67$. Leverage
rises as the cushion shrinks: the next loss hurts more.
\end{solution}

\begin{solution}{exo:m2:european-and-japanese-government-bonds:4}
First order: $0.5/21.8 = 229$ basis points. Second order: the loss $21.8\,x -
294.6\,x^2$ is largest at $x = 21.8/589 = 3.7\%$ and there equals $21.8^2/(2
\times 589) = 40.5\%$, below the 50\% cushion: the formula says the cushion is
never exhausted, which is false. Full repricing gives 351 basis points. The
quadratic understates losses far from the starting point because the true
price curve is flatter there than the parabola: convexity itself falls as
yields rise.
\end{solution}

\begin{solution}{exo:m2:european-and-japanese-government-bonds:5}
The problem was a dysfunction at the long end, forced selling by LDI funds,
not the stance of monetary policy; the Bank bought only what the funds were
selling, long conventional and index-linked \omterm{def:m2:european-and-japanese-government-bonds:ldi}{gilts}, for as long as it took the
funds to raise capital, and sold the \omterm{def:m2:european-and-japanese-government-bonds:ldi}{gilts} back afterwards. \omterm{def:m2:central-banks-and-the-short-rate:qe}{Quantitative easing}
buys across maturities, for years, to lower yields; the Bank was at the same
time raising rates and preparing to shrink its balance sheet, and a
temporary, targeted operation kept the two apart.
\end{solution}

\begin{solution}{exo:m2:european-and-japanese-government-bonds:6}
The market: the central bank buys whatever is offered at the target. If the
market expects the target to be raised or abandoned, holding bonds bought at
the target means a loss when the yield jumps, so holders sell to the central
bank first; its purchases, and its balance sheet, grow fastest exactly when
the policy is least credible, and it takes the loss when the target goes.
\end{solution}

\begin{solution}{exo:m2:european-and-japanese-government-bonds:7}
After $+100$: \omterm{def:m2:european-and-japanese-government-bonds:ldi}{gilts} worth 808.4 million, sales $2 \times 500 - 808.4 =
\text{GBP}~191.6$ million. After $+160$: \omterm{def:m2:european-and-japanese-government-bonds:ldi}{gilts} 715.6 million, sales GBP~284.4
million. To survive 250 basis points without selling, the cushion must cover
the loss at 250: 39.9\% of the holding (the fund's 50\% does).
\end{solution}

\begin{solution}{exo:m2:european-and-japanese-government-bonds:8}
Rising long yields reduce the present value of a defined-benefit scheme's
liabilities more than that of its equities and short bonds: the Bank's own
illustration notes that the scheme might be better off overall. The loss was
concentrated in the leveraged LDI fund, whose cushion was eroded, and the
crisis was one of liquidity: collateral had to be posted in hours, while the
schemes' other assets could not be turned into cash that fast.
\end{solution}

\begin{solution}{pb:m2:european-and-japanese-government-bonds:1}
\textbf{1.} Dirty price 63.41 per 100 (the \omterm{def:m2:european-and-japanese-government-bonds:ldi}{gilt} holds GBP~1.577 billion of face);
repo GBP~500 million; cushion GBP~500 million.
\textbf{2.} \omterm{def:m2:government-bonds:duration}{Modified duration} 21.84; the fund's DV01 is $10^9 \times 21.84 \times
10^{-4} = \text{GBP}~2.18$ million per basis point.
\textbf{3.} Leverage 2; 0.44\% of capital per basis point.
\textbf{4.} 1\% of the cushion is GBP~5 million: about 2.3 basis points.
\textbf{5.} To hedge more of its long liabilities than its bond holdings alone
would, while keeping growth assets: leverage buys duration without selling
equities.
\textbf{6.} \omterm{def:m2:european-and-japanese-government-bonds:ldi}{Gilts} GBP~808.4 million, cushion 308.4 million (62\% of the
original), leverage 2.62.
\textbf{7.} \omterm{def:m2:european-and-japanese-government-bonds:ldi}{Gilts} GBP~715.6 million, cushion 215.6 million (43\%), leverage 3.32.
\textbf{8.} GBP~191.6 million after 100; 284.4 million after 160.
\textbf{9.} At least GBP~50 billion against about 12 billion of daily volume in
those maturities, several days of the market's whole turnover in hours; the
price falls until other buyers appear, which raises yields and calls for more
sales.
\textbf{10.} LDI funds also held interest-rate and inflation swaps in which
they received fixed; rising rates made those swaps lose value, and
variation margin had to be posted in cash daily.
\textbf{11.} 351 basis points.
\textbf{12.} It sells the \omterm{def:m2:european-and-japanese-government-bonds:ldi}{gilts} it holds as collateral to recover its loan,
adding to the supply that caused the loss.
\textbf{13.} 39.9\% of the holding, GBP~399 million.
\textbf{14.} A 100-basis-point rise costs 19.2\% of the holding: a cushion of
about a fifth, less than half of what 250 requires.
\textbf{15.} The same cushion protects a short-duration fund against much larger
yield moves than a long one: resilience is a yield move, cushion divided by
duration, not a percentage.
\textbf{16.} Because the long end had stopped functioning: without a buyer
the spiral would have forced sales by banks of their collateral and disrupted
funding markets. The operation was sized and timed to the dysfunction, not to
the inflation outlook.
\textbf{17.} To buy time for the funds to raise capital and reduce leverage,
without making purchases a monetary-policy tool at a time of tightening.
\textbf{18.} Its liabilities fell with long yields, by more than its non-LDI
assets: its funding position was, if anything, better; what it lacked was
cash, fast.
\textbf{19.} Named result: \emph{the exhaustion yield} of the fund is a rise of 351
basis points; after the September move of 160 its cushion is 43\% of the
starting value, and restoring leverage would require selling GBP~284 million of
\omterm{def:m2:european-and-japanese-government-bonds:ldi}{gilts} into a falling market.
\textbf{20.} Leverage that had to be topped up with cash faster than cash could
be raised, in a market that could not absorb the sales.
\end{solution}

\begin{solution}{iq:m2:european-and-japanese-government-bonds:1}
Italy borrows in a currency it does not issue, so it can default in euros;
Germany is the benchmark of the safest credit and the most liquid market. The
spread pays for Italy's credit risk, the liquidity difference, and
\omterm{def:m2:european-and-japanese-government-bonds:spread}{redenomination risk}, the chance of repayment in a new, weaker currency; in a
crisis, also for the flight to \omterm{def:m2:european-and-japanese-government-bonds:bund}{Bunds}, which lowers the benchmark.
\iqlookfor{credit without a printing press, liquidity, redenomination.}
\end{solution}

\begin{solution}{iq:m2:european-and-japanese-government-bonds:2}
After the government's fiscal announcement of 23 September, long \omterm{def:m2:european-and-japanese-government-bonds:ldi}{gilt} yields
rose faster than any time since 2000, 160 basis points in a few days.
Pension-fund LDI strategies, leveraged through repo and swaps, faced
collateral calls they could meet fastest by selling \omterm{def:m2:european-and-japanese-government-bonds:ldi}{gilts}, which pushed yields
higher: a spiral. On 28 September the Bank of England bought long \omterm{def:m2:european-and-japanese-government-bonds:ldi}{gilts}
temporarily, 19.3 billion pounds by 14 October, and regulators then required
LDI funds to withstand a 250-basis-point move.
\iqlookfor{the trigger, the leverage mechanism, and the targeted backstop.}
\end{solution}

\begin{solution}{iq:m2:european-and-japanese-government-bonds:3}
The central bank targets a long yield and buys whatever is offered there; it
gives up control of its balance sheet. Exiting risks a jump in yields as the
buyer of last resort withdraws, losses for holders (and for the central bank
on its holdings), and spillovers to other markets whose investors had moved
abroad for yield. Exits are therefore done in steps, loosening the band
before ending the target.
\iqlookfor{price versus quantity control and the credibility problem.}
\end{solution}

\begin{solution}{iq:m2:european-and-japanese-government-bonds:4}
Decompose the change: issuer's yield change minus benchmark's. In a crisis the
benchmark often falls while the issuer rises. A hedge that shorts the
benchmark against a long in the issuer loses on the benchmark's rally; a
trader positioned for the spread must size both legs by DV01 and know which leg
is doing the moving.
\iqlookfor{the decomposition and its hedging consequence.}
\end{solution}

\begin{solution}{iq:m2:european-and-japanese-government-bonds:5}
OMT (2012): purchases without ex ante limit, conditional on an adjustment or
precautionary programme with the European Stability Mechanism. TPI (2022):
purchases against unwarranted, disorderly market dynamics that threaten the
transmission of policy, conditional on four criteria of fiscal and
macroeconomic soundness, without a programme. The first is a lender of last
resort with conditions; the second a transmission tool with eligibility.
\iqlookfor{the conditionality of each.}
\end{solution}

\begin{solution}{iq:m2:european-and-japanese-government-bonds:6}
Choose the shock from history at the relevant speed and horizon (how far
yields can move in the time it takes to raise capital), from the extremes of
comparable markets, and from reverse stress: the move that exhausts capital.
The test must capture full repricing rather than duration, liquidity (how much
can be sold, at what impact, in that time), margin and collateral calls on
derivatives, the counterparty's reaction (haircut increases), and the
feedback when many funds act alike.
\iqlookfor{reverse stress, liquidity and feedback, not only a larger number.}
\end{solution}
